\chapter{Pruebas y validación}
\label{ch:pruebas_validacion}

Este capítulo recoge únicamente evidencias verificadas durante el desarrollo de
Harmony Hub. No se incluyen resultados no comprobados ni afirmaciones sobre
comportamientos que no hayan quedado respaldados por ejecución real de la app,
por validaciones técnicas reproducibles o por métricas obtenidas del sistema.

\section{Plan de pruebas}
\label{sec:plan_pruebas}

La validación del proyecto se ha organizado en cuatro capas complementarias:
comprobaciones técnicas de consistencia estática sobre frontend y backend,
pruebas manuales del flujo funcional principal, ejecución de casos reales de
generación de playlist y revisión del estado del componente supervisado a
partir de sus metadatos, sus puertas de activación y el resumen observacional
por mood.

\begin{table}[H]
\centering
\scriptsize
\renewcommand{\arraystretch}{1.1}
\setlength{\tabcolsep}{3pt}
\begin{tabularx}{\textwidth}{|P{0.9cm}|P{2.2cm}|Y|P{2.4cm}|P{2.8cm}|P{1.9cm}|}
\hline
\textbf{ID} & \textbf{Precondición} & \textbf{Pasos} & \textbf{Respuesta esperada} & \textbf{Respuesta obtenida} & \textbf{Evidencia} \\
\hline
P01 & Código cliente actualizado y dependencias resueltas. & Ejecutar \texttt{dart analyze} sobre el proyecto Flutter. & Ausencia de errores de análisis en el cliente. & Superada. El análisis devolvió No issues found. & Salida de consola verificada \\
\hline
P02 & Backend actualizado y entorno Python disponible. & Ejecutar \texttt{python3 -m compileall app} en el backend. & Ausencia de errores de compilación sintáctica. & Superada. La compilación recorrió rutas, esquemas y servicios sin error. & Salida de consola verificada \\
\hline
P03 & Usuario autenticado y cuenta de Spotify conectada. & Ejecutar login, conexión con Spotify, check-in, recomendación, creación de playlist y feedback. & Flujo continuo sin bloqueo funcional. & Superada en validación manual. Todos los flujos principales quedaron operativos. & Observación manual del flujo completo \\
\hline
P04 & Sesión iniciada, Spotify autorizado y datos de check-in de relajación preparados. & Ejecutar el caso C1 con objetivo relajación, entorno quiet, preferencia instrumental e intensidad suave. & Recomendación coherente y playlist real creada en Spotify. & Superada. Se generó la playlist Harmony Hub · Relajacion y el flujo finalizó con feedback. & Logs de backend y caso C1 en Tabla~\ref{tab:casos_funcionales_playlist} \\
\hline
P05 & Sesión iniciada, Spotify autorizado y datos de check-in de foco preparados. & Ejecutar el caso C2 con objetivo foco, entorno medido, preferencia con voz e intensidad suave. & Recomendación orientada a concentración, activación del modelo de sesión y playlist real creada en Spotify. & Superada. Se generó la playlist Harmony Hub · Foco con trazas \texttt{ml\_enabled = True} y \texttt{use\_environment = True}. & Logs de backend y caso C2 en Tabla~\ref{tab:casos_funcionales_playlist} \\
\hline
P06 & Sesión iniciada, Spotify autorizado y datos de check-in de energía preparados. & Ejecutar el caso C3 con objetivo energía, usuario cansado y sin entorno acústico. & Playlist real creada manteniendo operatividad, incluso si el caso es difícil. & Superada. La playlist se generó y el sistema completó el flujo activando mecanismos de refuerzo sin bloquear la sesión. & Logs de backend y caso C3 en Tabla~\ref{tab:casos_funcionales_playlist} \\
\hline
P07 & Modelo entrenado, metadatos disponibles en el backend y configuración runtime cargada. & Revisar el archivo de metadatos del modelo supervisado y la configuración operativa del backend. & Umbrales explícitos, motivo de disponibilidad, resumen observacional por mood y umbral operativo de decisión. & Superada. Los metadatos exponen volumen de datos, cobertura mínima por clase y puerta global activa, mientras que la configuración runtime fija el umbral operativo de confianza. & Archivo \texttt{\seqsplit{session\_mode\_model\_metadata.json}} y configuración de backend \\
\hline
P08 & Backend actualizado y suite de pruebas disponible. & Ejecutar \texttt{python -m unittest discover -s tests} en el backend. & Todas las pruebas automáticas deben pasar. & Superada. La suite ejecutó 38 pruebas y finalizó con estado \texttt{OK}. & Salida de consola verificada \\
\hline
\end{tabularx}
\caption{Plan de pruebas ejecutado sobre Harmony Hub con precondiciones, pasos y evidencia.}
\label{tab:plan_pruebas}
\end{table}

\section{Pruebas funcionales}
\label{sec:pruebas_funcionales}

Las pruebas funcionales se han centrado en comprobar que el sistema cumple los
requisitos descritos en el Capítulo \ref{ch:analisis_problema}. La evidencia
obtenida confirma que el flujo principal de la aplicación se encuentra operativo
de extremo a extremo.

En validación manual se comprobó correctamente el acceso autenticado del usuario
(\textbf{RF1}), la navegación posterior a la conexión con Spotify
(\textbf{RF7}), la realización del check-in contextual (\textbf{RF3} y
\textbf{RF4}), la generación de una recomendación de sesión (\textbf{RF5}), la
trazabilidad mediante identificador de recomendación (\textbf{RF6}), la
creación de playlists reales en Spotify (\textbf{RF8}), el cierre con feedback
(\textbf{RF10}) y el acceso a historial, aprendizaje, modos y perfil
(\textbf{RF9}, \textbf{RF11} y \textbf{RF12}).

La comprobación funcional más relevante ha sido la ejecución de tres casos
completos de generación:

\medskip
\begingroup
\scriptsize
\renewcommand{\arraystretch}{1.1}
\setlength{\tabcolsep}{3pt}
\begin{longtable}{|P{1.0cm}|P{2.9cm}|P{2.1cm}|P{5.1cm}|P{1.8cm}|}
\caption{Casos funcionales reales ejecutados para la generación de playlists.}
\label{tab:casos_funcionales_playlist}\\
\hline
\textbf{Caso} & \textbf{Objetivo y contexto} & \textbf{Modo recomendado} & \textbf{Observaciones verificadas} & \textbf{Resultado} \\
\hline
\endfirsthead
\hline
\textbf{Caso} & \textbf{Objetivo y contexto} & \textbf{Modo recomendado} & \textbf{Observaciones verificadas} & \textbf{Resultado} \\
\hline
\endhead
\hline
\multicolumn{5}{r}{\textit{Continúa en la página siguiente}} \\
\endfoot
\hline
\endlastfoot
C1 & Relajación, usuario estresado, energía baja-media, entorno quiet, preferencia instrumental e intensidad suave. & Relajación / estresado / suave & Se generó playlist Harmony Hub · Relajacion. El backend construyó perfil con géneros semilla ambient, chill y classical; además cerró la sesión con feedback sin error funcional. & Correcto \\
\hline
C2 & Foco, usuario estresado, energía media, entorno medido con picos intermitentes, preferencia con voz e intensidad suave, ejecutado desde dispositivo Android físico. & Foco / estresado / media & Se activó el modelo supervisado a nivel de sesión, el entorno se incorporó a la recomendación y a la generación, y se creó la playlist Harmony Hub · Foco en Spotify. Las trazas mostraron \texttt{ml\_enabled = True}, \texttt{selection\_source = session\_ml} y \texttt{use\_environment = True}. & Correcto \\
\hline
C3 & Energía, usuario cansado, energía baja, sin uso de entorno, intensidad suave y perfil de popularidad mainstream. & Energía / cansado / suave & Se generó playlist Harmony Hub · Energia. El sistema completó el ranking principal sobre msd\_tracks y materializó la sesión final en Spotify sin bloqueo funcional. & Correcto \\
\end{longtable}
\endgroup

Estos tres casos permiten afirmar que el flujo de generación está implementado
y es operativo. Además, permiten comprobar que el comportamiento final no
depende de un único tipo de sesión, sino que se mantiene correcto en
relajación, foco y energía, con construcción de perfil, ranking sobre catálogo
MSD, materialización final en Spotify y activación efectiva de la capa
supervisada cuando la evidencia acumulada lo permite.

El caso C2 resulta especialmente relevante como evidencia de cierre del
proyecto, ya que reúne en una misma ejecución autenticación PKCE, lectura del
perfil real de Spotify, recomendación con modelo supervisado activo, uso
efectivo del entorno acústico medido desde el móvil, materialización final de
la playlist y mantenimiento automático posterior del conjunto supervisado.

En consecuencia, la validación funcional final permite sostener tres
afirmaciones que describen con precisión el estado actual del sistema. En
primer lugar, el entorno acústico ya no debe considerarse una variable
experimental o residual, sino una señal realmente utilizada por la generación
musical cuando se activa desde el cliente. En segundo lugar, las entradas del
usuario están acotadas, validadas y trazadas a lo largo de todo el flujo, desde
el check-in hasta la playlist y el feedback. En tercer lugar, la aplicación
ejecuta ya un recorrido completo y reproducible de extremo a extremo, con
persistencia documental, trazabilidad entre recomendación y playlist, y
retroalimentación posterior del aprendizaje.

\subsection*{Cobertura de los casos de uso}

La revisión funcional también permite contrastar si los casos de uso descritos
en el Capítulo~\ref{ch:analisis_problema} han quedado realmente cubiertos por
la aplicación y por su backend asociado.

\begingroup
\small
\renewcommand{\arraystretch}{1.1}
\setlength{\tabcolsep}{4pt}
\begin{longtable}{|P{1.0cm}|P{4.2cm}|P{6.0cm}|P{1.8cm}|}
\caption{Comprobación de cobertura de los casos de uso principales.}
\label{tab:cobertura_casos_uso}\\
\hline
\textbf{ID} & \textbf{Caso de uso} & \textbf{Evidencia verificada} & \textbf{Estado} \\
\hline
\endfirsthead
\hline
\textbf{ID} & \textbf{Caso de uso} & \textbf{Evidencia verificada} & \textbf{Estado} \\
\hline
\endhead
\hline
\multicolumn{4}{r}{\textit{Continúa en la página siguiente}} \\
\endfoot
\hline
\endlastfoot
CU1 & Registrarse e iniciar sesión & Flujo disponible en cliente mediante Firebase Authentication y validado durante las pruebas manuales de acceso. & Cumplido \\
\hline
CU2 & Recuperar contraseña & Implementado mediante envío de correo de restablecimiento desde \texttt{AuthService} y accesible desde la pantalla de autenticación. & Cumplido \\
\hline
CU3 & Conectar cuenta con Spotify & Flujo OAuth con PKCE operativo a través de \texttt{/spotify/login-url}, \texttt{/spotify/exchange} y \texttt{/spotify/me}. & Cumplido \\
\hline
CU4 & Realizar una sesión musical contextual & Comprobado en los tres casos funcionales ejecutados, con check-in, recomendación, playlist y trazabilidad. & Cumplido \\
\hline
CU5 & Consultar historial personal & Disponible en la app mediante \texttt{history\_firestore\_service} y validado sobre colecciones de check-ins, playlists y feedback. & Cumplido \\
\hline
CU6 & Guardar un modo rápido & Implementado sobre la colección \texttt{saved\_preset\_modes} y disponible en el módulo de presets del cliente. & Cumplido \\
\hline
CU7 & Enviar feedback sobre una sesión & Validado en el cierre de sesión, con persistencia pública y privada y efecto posterior sobre el aprendizaje. & Cumplido \\
\hline
CU8 & Consultar aprendizaje personal & Cubierto por la pantalla de aprendizaje y por los documentos agregados en \texttt{user\_generation\_preferences}. & Cumplido \\
\end{longtable}
\endgroup

\section{Pruebas no funcionales}
\label{sec:pruebas_no_funcionales}

Las pruebas no funcionales se han abordado desde cinco perspectivas: calidad
técnica del código, robustez frente a dependencias externas, trazabilidad de la
sesión, seguridad del dato y preparación del componente de aprendizaje.

\begin{table}[H]
\centering
\small
\renewcommand{\arraystretch}{1.1}
\begin{tabularx}{\textwidth}{|P{1.4cm}|Y|P{2.8cm}|P{3.1cm}|}
\hline
\textbf{ID} & \textbf{Comprobación} & \textbf{Requisito asociado} & \textbf{Resultado verificado} \\
\hline
NF1 & Análisis estático del cliente Flutter. & RNF1, RNF6 & Superado. \texttt{dart analyze} no devolvió incidencias. \\
\hline
NF2 & Compilación completa del backend Python. & RNF6, RNF7 & Superado. \texttt{compileall} recorrió rutas, servicios, esquemas y scripts sin error. \\
\hline
NF3 & Continuidad del sistema cuando Spotify limita búsquedas o reduce respuestas externas. & RNF5, RNF8 & Superado. El sistema mantiene el ranking sobre catálogo MSD, limita reintentos excesivos y preserva la trazabilidad del flujo. \\
\hline
NF4 & Persistencia protegida de datos emocionales en colecciones públicas y privadas. & RNF3, RNF4 & Superado. La arquitectura separa bloque público y bloque cifrado en Firestore. \\
\hline
NF5 & Disponibilidad condicionada del modelo supervisado por calidad, volumen de datos y umbral mínimo de confianza. & RNF7, RNF8, RNF9 & Superado. Existe control explícito mediante data gate, quality gate, umbral operativo de decisión y mantenimiento automático del modelo tras el feedback. \\
\hline
\end{tabularx}
\caption{Resumen de pruebas no funcionales verificadas.}
\label{tab:pruebas_no_funcionales}
\end{table}

La evidencia reunida combina análisis estático, compilación satisfactoria,
suite automática de backend, trazas de ejecución y validación manual de flujos
completos sobre la aplicación móvil. Con ello se ha cubierto de forma
suficiente la validación funcional y técnica requerida para la entrega final.

\section{Correcciones realizadas tras las pruebas}
\label{sec:correcciones_pruebas}

Las pruebas no solo han servido para confirmar comportamientos correctos, sino
también para localizar y corregir fallos concretos del sistema.

Una de las incidencias detectadas apareció en el cierre de feedback desde el
cliente móvil. La app intentaba escribir en la colección \texttt{feedback} con
un documento que no coincidía exactamente con las reglas activas de Firestore,
lo que provocaba un error de permisos. La corrección consistió en alinear el
payload enviado por el cliente con el esquema realmente permitido por
\texttt{firestore.rules}, incluyendo los campos obligatorios y el nombre exacto
del indicador de bloque cifrado.

Otra corrección relevante ha afectado al flujo de generación de playlist cuando
las búsquedas externas entran en limitación de tasa. En lugar de quedar bloqueado
durante esperas muy largas, el sistema acota reintentos, conserva el ranking
musical construido sobre el catálogo MSD y mantiene un flujo de materialización
controlado cuando es viable. Esto resulta coherente con la arquitectura final:
Spotify ya no actúa como fuente principal de descubrimiento dinámico, sino como
capa de autorización, contexto musical básico y materialización final.

La validación del componente aprendido llevó a refinar la política
de activación del modelo. La versión actual no se basa en desbloquear el
aprendizaje tras un número fijo de sesiones, sino en exigir evidencia
suficiente a nivel global, junto con un umbral operativo de confianza en la
sesión concreta. La observabilidad por \textit{mood} se conserva como señal de
auditoría y lectura del dataset, pero ya no actúa como puerta dura
independiente del clasificador global. Esto permitió alinear mejor el
comportamiento del sistema con una explicación académica más prudente y más
fiel al estado real de los datos.

En esa misma línea se corrigió el comportamiento binario del aprendizaje por
mood. Antes, si el estado emocional actual no superaba su puerta local,
los pesos de aprendizaje podían caer a cero incluso cuando ya existía memoria
global del usuario. La implementación actual sustituye ese salto rígido por una
aplicación progresiva, apoyada en
\texttt{mood\_learning\_application\_factor}, lo que permite representar mejor
la transición entre contexto puro, aprendizaje parcial y perfil estable.

También se automatizó el mantenimiento del modelo supervisado. El cierre del
feedback dispara ahora una auditoría en segundo plano que compara el número de
ejemplos etiquetados con los metadatos previos y decide si debe reentrenar o
si puede omitir la operación. Esto eliminó la dependencia de comandos manuales
para mantener el artefacto alineado con el conjunto real de
\texttt{training\_session\_examples}.

Una corrección adicional, relevante para la coherencia funcional del sistema,
afectó al enlace entre la recomendación mostrada y la playlist realmente
generada. En una versión intermedia del cliente, la llamada de generación podía
arrastrar la preferencia de intensidad original del check-in en lugar de la
intensidad final contenida en \texttt{recommended\_mode}. La versión actual ya
alinea ambos pasos, de forma que la playlist se construye a partir del modo de
sesión realmente recomendado y no de una preferencia previa potencialmente
distinta.

También se corrigió una incoherencia en el historial expuesto desde backend.
El endpoint \texttt{/history/me} inicialmente reconstruía una vista parcial
basada en estructuras en memoria, lo que podía divergir del estado persistido
real de Firestore. La implementación actual lo alinea con la fuente documental
persistida del usuario y deja el descifrado de los bloques privados en la capa
cliente, que es donde reside la clave local.

Durante la validación también se revisaron elementos de interfaz y
persistencia que bloqueaban la experiencia, como desajustes de navegación,
problemas de overflow visual y la correspondencia entre reglas de
Firestore y documentos generados desde el cliente.

\section{Caso de prueba longitudinal del componente de aprendizaje}
\label{sec:caso_prueba_ml_evidencia}

Además de validar la operatividad general de la aplicación, resultaba
necesario comprobar que el componente de aprendizaje pudiera observarse como un
proceso acumulativo en el tiempo y no solo como una afirmación teórica. Para
ello se contrastaron trazas de inferencia, crecimiento del conjunto de
entrenamiento y metadatos persistidos de un ciclo real de mantenimiento sobre
la colección \texttt{training\_session\_examples}.

\begin{table}[H]
\centering
\small
\renewcommand{\arraystretch}{1.1}
\begin{tabularx}{\textwidth}{|P{1.6cm}|P{3.2cm}|Y|P{2.2cm}|}
\hline
\textbf{Fase} & \textbf{Comprobación} & \textbf{Evidencia observada} & \textbf{Resultado} \\
\hline
\textbf{T1} & Mantenimiento automático tras el feedback & El backend deja trazas \texttt{[ML MAINTENANCE] start}, \texttt{train} o \texttt{skip}, y persiste un artefacto de auditoría con fecha, ejemplos actuales y motivo de reentrenado u omisión. & Correcto \\
\hline
\textbf{T2} & Estado real del conjunto supervisado & La colección \texttt{training\_session\_examples} alcanzó \texttt{66} ejemplos etiquetados, con balance \texttt{54} útiles frente a \texttt{12} no útiles. & Correcto \\
\hline
\textbf{T3} & Disponibilidad del modelo global & Los metadatos indican \texttt{model\_gate\_passed = true}, \texttt{data\_gate\_passed = true} y \texttt{quality\_gate\_passed = true}, mientras que la configuración operativa del backend fija un umbral \texttt{min\_selected\_mode\_probability = 0.14}; en conjunto, esto confirma la activación del clasificador de sesión en ejecución real. & Correcto \\
\hline
\textbf{T4} & Aprendizaje individual observable & En la capa de aprendizaje por perfil, el sistema ya muestra aplicación completa (\texttt{mood\_learning\_application\_factor = 1.0}) y modo \texttt{stable\_weighted} para los bloques maduros de triste, estresado, cansado, neutral y feliz. & Correcto \\
\hline
\textbf{T5} & Activación integrada del flujo final & Una sesión real de foco/estresado ejecutada desde móvil físico dejó trazas \texttt{ml\_enabled = True}, \texttt{selection\_source = session\_ml}, \texttt{use\_environment = True} y finalizó con playlist creada correctamente en Spotify. & Correcto \\
\hline
\end{tabularx}
\caption{Caso de prueba longitudinal del componente de aprendizaje con evidencia real de mantenimiento, activación supervisada y aprendizaje individual.}
\label{tab:caso_prueba_ml_longitudinal}
\end{table}

La tabla anterior permite verificar tres hechos relevantes. En primer lugar, el
aprendizaje supervisado global ya puede activarse con evidencia operativa real
y trazable. En segundo lugar, el mantenimiento automático del backend hace que
esa activación no dependa de una intervención manual, sino de una auditoría
reproducible tras el feedback. En tercer lugar, la activación del clasificador
global convive con la capa individual del usuario, donde siguen apareciendo
modos como \texttt{stable\_weighted} y factores de aplicación completa para los
cinco moods observados en el perfil único.


\subsection*{Evidencia directa desde los artefactos JSON del modelo}

Además de las trazas de consola, el repositorio conserva tres artefactos JSON
que permiten demostrar de forma directa qué ha aprendido el clasificador y cómo
varía su estado al cerrarse una sesión nueva con feedback.

\begin{table}[H]
\centering
\small
\renewcommand{\arraystretch}{1.1}
\begin{tabularx}{\textwidth}{|P{2.7cm}|P{5.2cm}|Y|}
\hline
\textbf{Artefacto JSON} & \textbf{Campos observados} & \textbf{Lectura evidenciable} \\
\hline
\textbf{Audit del modelo} & \texttt{current\_labeled\_examples = 66}, \texttt{previous\_metadata\_examples = 65}, \texttt{trained = true}. & Demuestra variación real tras una sesión nueva: el conjunto etiquetado crece en una unidad y el backend ejecuta reentrenamiento automático en lugar de limitarse a conservar un modelo estático. \\
\hline
\textbf{Metadatos del modelo} & \texttt{total\_examples = 66}, \texttt{class\_counts = \{54,12\}}, \texttt{data\_gate\_passed = true}, \texttt{quality\_gate\_passed = true}, \texttt{model\_gate\_passed = true}, \texttt{quality\_score = 62.21}. & Confirma el estado aprendido final del clasificador: ya no solo existe, sino que supera las puertas de datos y calidad y queda disponible para intervenir en ejecución real. \\
\hline
\textbf{Model card} & Coeficientes positivos y negativos más influyentes, por ejemplo \texttt{use\_environment = +1.0925}, \texttt{recommended\_mode = relajacion\_estresado\_suave = +0.6561} y \texttt{stress\_level = -0.4719}. & Hace visible qué patrones ha internalizado el modelo y permite defender que el aprendizaje no es opaco: se puede inspeccionar qué rasgos empujan o frenan la probabilidad de utilidad estimada. \\
\hline
\end{tabularx}
\caption{Evidencia directa obtenida de los artefactos JSON persistidos por el componente supervisado.}
\label{tab:evidencia_json_modelo}
\end{table}

La primera fila es especialmente valiosa porque muestra una variación temporal
explícita. El archivo \texttt{\seqsplit{session\_mode\_model\_audit.json}}
indica que, antes del último mantenimiento automático, el metadato consolidado
del modelo reflejaba \texttt{65} ejemplos, mientras que después del cierre de
la nueva sesión pasa a \texttt{66}. No se trata, por tanto, de una fotografía
aislada del sistema, sino de una evidencia de actualización efectiva tras el
feedback.

La segunda fila permite ir más allá del simple contador. En
\texttt{\seqsplit{session\_mode\_model\_metadata.json}} también aparecen las
métricas \texttt{balanced\_accuracy\_mean = 0.5658},
\texttt{roc\_auc\_mean = 0.5818} y \texttt{f1\_mean = 0.8194}, junto con un
\texttt{quality\_score = 62.21}. Esto demuestra que el estado aprendido no solo
crece en volumen, sino que queda reevaluado y respaldado por métricas que
superan los mínimos operativos del proyecto.

La tercera fila aporta una evidencia cualitativa complementaria. En lugar de
limitarse a declarar que el modelo ``aprende'', el archivo
\texttt{\seqsplit{session\_mode\_model\_card.json}} deja ver qué ha aprendido.
Que \texttt{use\_environment} aparezca entre las señales positivas más fuertes
resulta coherente con la arquitectura contextual de Harmony Hub; que
\texttt{stress\_level} figure entre las señales negativas importantes también
encaja con el hecho de que algunos modos dejan de ser útiles cuando el estrés
sube demasiado. Este tipo de lectura hace al componente supervisado mucho más
defendible desde el punto de vista académico.


\section{Síntesis final de la validación del aprendizaje}
\label{sec:sintesis_final_aprendizaje}

La evidencia acumulada permite cerrar la validación con una lectura clara del
estado actual del sistema. Harmony Hub ya muestra aprendizaje individual del
perfil plenamente observable en los cinco estados emocionales trabajados
(triste, estresado, cansado, neutral y
feliz), con trazas longitudinales que permiten ver tanto la maduración
gradual de algunos bloques como la transición desde fases prudentes de
aplicación parcial hasta modos estables con factor completo.

Además, la implementación ya no se limita a conservar el clasificador
supervisado como componente integrado pero inactivo. El conjunto actual, con 66
ejemplos etiquetados, balance suficiente entre clases y cobertura completa por
mood, permite activar el modelo global en ejecución real y observar cómo
participa en la selección del modo de sesión cuando la confianza supera el
umbral operativo configurado. Por tanto, la validación final no concluye ya
solo que el sistema aprende, sino que demuestra también que la capa
supervisada interviene de manera trazable y compatible con el carácter híbrido
de la arquitectura.

\section{Resultados y discusión}
\label{sec:resultados_discusion}

La validación realizada permite sostener, en primer lugar, que Harmony Hub ha
alcanzado un grado de operatividad suficiente como sistema funcional completo.
El flujo formado por autenticación, check-in, recomendación, generación de
playlist, persistencia, historial y cierre con feedback se ha ejecutado de
extremo a extremo sobre la aplicación móvil y el backend real, con evidencia
reproducible tanto en la interfaz como en las trazas técnicas del sistema.

Esa misma validación permite defender además que el proyecto ya cumple tres
propiedades que al inicio solo podían plantearse como objetivo de diseño. La
primera es que el entorno acústico funciona de forma efectiva dentro del
pipeline de personalización, ya que se ha observado en ejecución real con uso
activado, fuerza de influencia positiva y consultas auxiliares asociadas al
contexto detectado. La segunda es que los inputs del usuario se mantienen
controlados y consistentes, al recogerse mediante opciones cerradas y al quedar
validados y persistidos en estructuras documentales alineadas entre cliente y
backend. La tercera es que el flujo completo del sistema resulta trazable de
extremo a extremo, porque cada sesión deja enlazados su check-in, su
recomendación, su playlist generada, su feedback y su efecto posterior sobre el
aprendizaje.

En segundo lugar, los resultados confirman que la capa de aprendizaje no debe
interpretarse como un bloque único, sino como dos mecanismos complementarios.
El aprendizaje individual del perfil musical ofrece una adaptación observable y
estable en los cinco estados emocionales principales del proyecto, mientras que
el clasificador supervisado global ya puede entrar en ejecución real cuando la
sesión presenta condiciones compatibles con su puerta de calidad y con su
umbral operativo de confianza. Esta distinción aporta valor metodológico porque
permite mostrar, dentro de una misma solución, tanto personalización
progresiva como intervención supervisada prudente.

En tercer lugar, la discusión de resultados pone de manifiesto una consecuencia
relevante del diseño adoptado: incluso con el modelo global activo, la utilidad
del sistema no depende de que toda la decisión recaiga sobre el clasificador.
La arquitectura híbrida basada en catálogo estructurado, reglas contextuales,
perfil aprendido y materialización real en Spotify mantiene un comportamiento
explicable y estable. El modelo supervisado añade capacidad de reordenación,
pero sigue operando dentro de un marco trazable de umbrales, deltas y
abstenciones locales.

Por último, la validación también delimita con claridad los márgenes de mejora.
La principal restricción ya no se encuentra en la implementación del flujo
funcional ni en la ausencia de evidencia longitudinal, sino en la posibilidad de
seguir refinando la separación estadística entre clases y la calibración del
modelo conforme crezcan los datos supervisados. A partir de esta observación,
el proyecto puede considerarse cerrado como sistema híbrido entregable y, al
mismo tiempo, abierto como base experimentable para futuras iteraciones
centradas en fortalecer todavía más el componente global.
